\chapter{Réalisation et implémentation}

\section{Introduction}

Le présent chapitre décrit la réalisation concrète de Strategic Copilot, plateforme SaaS d'aide à la décision pour la gestion des talents et des projets, telle qu'elle est livrée dans le dépôt \texttt{frontend-ui}. Il fait suite au chapitre de conception : les architectures, workflows et contrats y ont été formalisés ; il s'agit ici d'expliquer comment elles ont été implémentées, dans quel environnement de développement, et quelles interfaces sont effectivement utilisables par les profils manager, RH et talent.

La réalisation s'organise autour d'une application React (\texttt{copilot-ui}) consommant des webhooks n8n hébergés sur \texttt{n8nprod.aphelionxinnovations.com}, avec persistance PostgreSQL (Supabase) et enrichissement IA orchestré côté serveur (paramètre \texttt{use\_ai: true}, API compatibles OpenAI). Aucune fonctionnalité non présente dans le code source n'est décrite : chaque affirmation de ce chapitre est vérifiable dans les fichiers TypeScript, de configuration Vite/Vercel et dans le script d'intégration \texttt{n8n/route-validate-claims-merger-FIX.js}.

\section{Environnement de développement}

\subsection{Outils utilisés}

Le développement a été mené avec Visual Studio Code ou Cursor, Git pour le versionnement du dépôt \texttt{frontend-ui}, et npm pour la gestion des dépendances à la racine du monorepo. La compilation du présent rapport utilise LaTeX Workshop (MiKTeX) sur le dossier \texttt{rapport/}. Les tests manuels s'effectuent dans un navigateur moderne (Chrome ou Edge) contre l'instance n8n de production ou de préproduction, relayée en local par le proxy Vite.

Les scripts npm principaux, définis dans le \texttt{package.json} racine, sont \texttt{npm run dev} (serveur Vite sur le port 5173), \texttt{npm run build} (bundle de production dans \texttt{copilot-ui/dist}), \texttt{npm run typecheck} et \texttt{npm run build:check} (vérification TypeScript avant build). Le sous-dossier \texttt{copilot-ui} contient un \texttt{package.json} minimal qui délègue ces commandes à la racine, ce qui évite la duplication des dépendances React.

\subsection{Technologies retenues}

Le tableau~\ref{tab:stack-ch5} synthétise la stack effectivement déclarée dans les manifestes du projet.

\begin{table}[htbp]
\centering
\caption{Stack technique de Strategic Copilot (versions du \texttt{package.json} racine)}
\label{tab:stack-ch5}
\begin{tabular}{@{}ll@{}}
\toprule
\textbf{Composant} & \textbf{Version / bibliothèque} \\
\midrule
Runtime UI & React 19.2, React DOM 19.2 \\
Build & Vite 7.3, plugin \texttt{@vitejs/plugin-react-swc} \\
Langage & TypeScript 5.9 \\
Routage & React Router 7.13 \\
Données serveur & TanStack React Query 5.90 \\
HTTP & axios 1.13 \\
Styles & Tailwind CSS 4.1, plugins typography et react-aria \\
Formulaires & react-hook-form 7.54, zod 3.24 \\
i18n & i18next 25.8, react-i18next 16.5 \\
Graphiques & Recharts 2.15 \\
Composants & React Aria Components 1.15, Untitled UI \\
\bottomrule
\end{tabular}
\end{table}

Le backend métier n'est pas un serveur Node dédié dans ce dépôt : il est constitué des workflows n8n et de PostgreSQL. L'IA générative est invoquée dans les graphes n8n lorsque le corps de requête contient \texttt{use\_ai: true}, sans SDK OpenAI embarqué dans le bundle client.

\subsection{Configuration de l'environnement}

Le fichier \texttt{copilot-ui/.env} documente les variables essentielles. \texttt{VITE\_N8N\_BASE\_URL} pointe vers l'hôte n8n ; en développement, le client HTTP utilise des chemins relatifs \texttt{/webhook/...} pour profiter du proxy défini dans \texttt{vite.config.ts}. \texttt{VITE\_API\_BASE\_URL} peut surcharger la base des webhooks versionnés (\texttt{wmp-*}, \texttt{wmt-*}). Les clés Supabase (\texttt{VITE\_SUPABASE\_URL}, \texttt{VITE\_SUPABASE\_ANON\_KEY}) sont réservées au stockage fichier ; les requêtes analytiques passent par n8n.

Le fichier \texttt{copilot-ui/.env.production} fixe les chemins d'authentification et RH pour le déploiement Vercel (\texttt{VITE\_API\_LOGIN=/webhook/login}, \texttt{VITE\_API\_RH\_USERS\_LIST=/webhook/rh/users}, etc.). Le fichier \texttt{vercel.json} à la racine relaie \texttt{/webhook/:path*} vers \texttt{n8nprod.aphelionxinnovations.com}, ce qui aligne production et développement sur le même schéma d'URL côté navigateur.

La figure~\ref{fig:archi-dev-ch5} schématise l'environnement de développement local.

\begin{figure}[htbp]
\centering
\fbox{\parbox{0.92\textwidth}{\centering\vspace{0.8cm}
\textbf{Schéma d'exécution en développement}\\[0.4cm]
Navigateur $\rightarrow$ \texttt{localhost:5173} (Vite)\\
$\rightarrow$ proxy \texttt{/webhook}, \texttt{/api}, \texttt{/n8n-webhook}\\
$\rightarrow$ \texttt{n8nprod.aphelionxinnovations.com}\\
$\rightarrow$ PostgreSQL / OpenAI (nœuds n8n)\vspace{0.8cm}}}
\caption{Enchaînement des appels en mode développement (proxy Vite)}
\label{fig:archi-dev-ch5}
\end{figure}

\section{Implémentation du frontend}

\subsection{Structure du projet React/Vite}

La racine Vite est \texttt{copilot-ui} (\texttt{root: appRoot} dans \texttt{vite.config.ts}). Le point d'entrée \texttt{src/main.tsx} monte l'arbre React avec les providers globaux (\texttt{ThemeProvider}, \texttt{AuthProvider}, \texttt{QueryClientProviderWrapper}, \texttt{ToastProvider}, \texttt{CopilotProvider}, \texttt{WhatIfProvider}) et le routeur. L'alias \texttt{@} résout vers \texttt{src/}, ce qui simplifie les imports métier.

L'arborescence suit une séparation par responsabilité. Le répertoire \texttt{pages/} contient les écrans par workspace (\texttt{manager/}, \texttt{workspace/rh/}, \texttt{workspace/talent-workspace-pages}). Le répertoire \texttt{components/} regroupe les blocs réutilisables (Copilot, Mission Control, formulaires RH). Le répertoire \texttt{api/} centralise les appels HTTP typés ; \texttt{hooks/} encapsule React Query ; \texttt{features/manager/project-detail/} concentre la logique du modal Mission Control (traces d'agents, recompute). Le répertoire \texttt{types/} définit les contrats (\texttt{api.types.ts}, \texttt{crud-domain.ts}, \texttt{copilot.ts}).

Le principe fondamental du domaine est codifié en tête de \texttt{crud-domain.ts} : le front-end n'effectue aucun recalcul de scores ni de recommandations. L'extrait suivant en fixe la portée.

\begin{lstlisting}[caption={Principe pas de calcul front (\texttt{crud-domain.ts})},label={lst:crud-domain}]
/** Champs analytiques : lecture seule depuis l'API.
 *  Le frontend ne recalcule pas scores ni recommandations. */
export interface Project {
  viability_score?: number | null;
  risk_score?: number | null;
  decision?: string | null;
}
\end{lstlisting}

\subsection{Gestion du routage}

React Router 7 déclare trois layouts protégés par rôle dans \texttt{main.tsx}. \texttt{ProtectedRoute} redirige vers \texttt{/login} si la session est absente et vérifie le tableau \texttt{roles} pour les routes sensibles. Les chemins legacy (\texttt{/projects/:id}, \texttt{/decision-log}) redirigent vers les workspaces via \texttt{LegacyProjectDetailRedirect} et composants associés, préservant les liens profonds.

L'espace manager vit sous \texttt{/workspace/manager/*} avec \texttt{ManagerWorkspaceLayout} ; l'espace RH sous \texttt{/workspace/rh/*} avec \texttt{RhWorkspaceLayout} ; l'espace talent sous \texttt{/workspace/talent/*} avec \texttt{TalentWorkspaceAppLayout}. La route \texttt{/workspace/rh/manager-requests} est partagée entre manager et RH (\texttt{RhManagerRequestsEntry}) pour le traitement des demandes. Le panneau \texttt{CopilotPanel} est rendu au niveau racine du routeur, visible sur les écrans autorisés selon le contexte fourni par \texttt{CopilotProvider}.

\subsection{Gestion des états avec React Query}

\texttt{QueryClientProviderWrapper} configure un \texttt{QueryClient} partagé : \texttt{staleTime} de 60 secondes, \texttt{gcTime} de cinq minutes, deux tentatives maximum sur erreur transitoire, \texttt{refetchOnWindowFocus} activé. Chaque domaine métier possède des hooks dédiés : \texttt{useDashboard}, \texttt{use-portfolio}, \texttt{useProjects}, \texttt{use-manager-risk-data}, \texttt{use-rh-actions-query}, \texttt{use-talent-workspace-page-query}, \texttt{use-copilot-query}.

Les mutations invalident sélectivement les caches. Après une affectation projet, \texttt{useProjects} peut déclencher \texttt{/webhook/api/copilot/recompute} si \texttt{VITE\_TRIGGER\_PROJECT\_RECOMPUTE\_AFTER\_ASSIGN=1}. Après exécution Strategist, les requêtes risques et portefeuille sont rafraîchies. \texttt{CrossRoleQuerySync} propage les invalidations lors d'un changement de session pour éviter l'affichage de données d'un autre rôle.

L'état UI local (onglets Mission Control, modales, formulaires) reste dans les composants React ou dans des contextes (\texttt{WhatIfProvider} pour le simulateur global). Cette séparation évite de mélanger cache serveur et brouillons utilisateur.

\subsection{Design system et Tailwind CSS}

Tailwind CSS 4 est intégré via \texttt{@tailwindcss/vite}. Les styles globaux (\texttt{globals.css}) définissent les variables de thème clair/sombre et les surcharges Recharts pour la lisibilité des graphiques. Les composants s'appuient sur \texttt{tailwind-merge} et sur la bibliothèque Untitled UI (icônes, boutons, modales). React Aria assure l'accessibilité clavier et les rôles ARIA sur les champs et tableaux.

L'internationalisation couvre le workspace manager en français, anglais et arabe (\texttt{manager-workspace-locales.ts}, chargé depuis \texttt{i18n.ts}). Les dates relatives utilisent \texttt{localeForDateFormatting} (\texttt{ui-locale.ts}) pour adapter le formatage au locale actif. Le \texttt{ThemeProvider} permet le basculement clair/sombre sans rechargement.

\subsection{Workspaces Manager, RH et Talent}

L'espace manager regroupe le pilotage opérationnel : \texttt{DashboardPage} consomme \texttt{useDashboard} et les blocs \texttt{matchmaker}/\texttt{analyst} ; \texttt{ProjectsPageManager} ouvre \texttt{ProjectMissionControlModal} ; \texttt{TeamPage} et \texttt{TalentDetailPage} appellent les webhooks équipe et Watchdog ; \texttt{RisksPage} agrège \texttt{/webhook/api/project/risks} ; \texttt{ManagerRhRequestsPage} et \texttt{TalentRequestsPage} gèrent les escalades ; \texttt{ReportsPage} et \texttt{DecisionLogPage} complètent la gouvernance.

L'espace RH implémente le référentiel et la politique RH : pages \texttt{rh-dashboard-page}, \texttt{rh-employees-page}, \texttt{rh-skills-catalog-page}, \texttt{rh-critical-gaps-page}, \texttt{rh-training-plans-page}, \texttt{rh-mobility-page}, \texttt{rh-org-alerts-page}, plus \texttt{rh-accounts-workspace-page} et \texttt{rh-sessions-workspace-page} branchés sur \texttt{rhService.ts} (\texttt{/webhook/rh/users}, \texttt{/webhook/rh/sessions}).

L'espace talent expose huit vues alimentées par \texttt{fetchTalentWorkspacePage} sur \texttt{/api/workspace/talent/dashboard}, \texttt{.../projects}, \texttt{.../tasks}, \texttt{.../workload}, \texttt{.../skills}, \texttt{.../trainings}, \texttt{.../notifications} et \texttt{.../profile}. La mise à jour d'une tâche passe par \texttt{talent-task-update.api.ts} (\texttt{/api/talent/tasks/:id}, méthode configurable).

\section{Implémentation du backend}

\subsection{Architecture générale des workflows n8n}

Les workflows n8n ne sont pas exportés dans le dépôt front-end ; leur contrat est inféré des chemins HTTP et des noms \texttt{WF\_*} référencés dans le code. Chaque flux suit généralement le patron Webhook $\rightarrow$ validation JWT $\rightarrow$ nœuds Postgres $\rightarrow$ branche \texttt{use\_ai} $\rightarrow$ Respond to Webhook. Le script \texttt{route-validate-claims-merger-FIX.js} montre que la fusion des claims et le routage manager team sont traités en amont des handlers.

En développement, trois préfixes coexistent : \texttt{/webhook} (chemins bruts), \texttt{/api} (réécrit en \texttt{/webhook/api} sauf \texttt{/api/rh/actions} vers le webhook UUID \texttt{c8bae94d-...}), et \texttt{/n8n-webhook} (réécrit vers \texttt{/webhook} pour les tâches \texttt{wmt-*-v1}).

\subsection{Principaux workflows réellement implémentés}

La réalisation couvre plusieurs familles de workflows, toutes consommées par \texttt{copilot-ui}.

La famille \textbf{manager projet} (\texttt{WF\_Manager\_Projects}, \texttt{wmp-detail-v1}, \texttt{wmp-update-v1}, \texttt{wmp-assign-v1}, \texttt{wmp-unassign-v1}) assure liste, détail, mise à jour et affectations. La famille \textbf{tâches} (\texttt{WF\_Manager\_Project\_Tasks}, \texttt{wmt-list-v1} à \texttt{wmt-delete-v1}) est implémentée dans \texttt{projectTasksApi.ts}. La famille \textbf{équipe} (\texttt{WF\_Manager\_Team}, \texttt{wmt-detail-v1/manager/team/:talentId}) alimente \texttt{TalentDetailPage}.

La famille \textbf{agentique} comprend \texttt{WF\_Project\_Analysis} (\texttt{POST /webhook/api/project/details}), \texttt{WF\_Project\_Risk} (\texttt{.../risks}), \texttt{WF\_Talent\_Matching} (\texttt{.../talents}), \texttt{WF\_Project\_Viability} (\texttt{.../viability}), le what-if (\texttt{.../what-if}), le recompute (\texttt{.../copilot/recompute}), le Strategist (\texttt{.../strategist/propose|execute}) et le Helper (\texttt{.../helper/chat}, \texttt{.../validations}). Les traces affichées dans Mission Control proviennent de \texttt{buildAgentTraces} dans \texttt{agent-trace.ts}.

La famille \textbf{RH et demandes} regroupe \texttt{WF\_Manager\_RH\_Actions}, les agrégats \texttt{/api/rh/*} et le PATCH dédié sur \texttt{/api/rh/actions/:id}. La famille \textbf{rapports} expose \texttt{/webhook/reports/generate-board-pack}, \texttt{generate-project-dossier}, \texttt{history}, \texttt{send-email} et \texttt{schedule}. La famille \textbf{auth} couvre login, refresh, logout et profil.

\subsection{Intégration des services OpenAI}

Aucune clé OpenAI n'est présente dans le code React. L'intégration est réalisée dans n8n : les appels front-end passent \texttt{use\_ai: true} lorsque l'analyse enrichie est requise, par exemple dans \texttt{agents.api.ts} pour les risques et le matching, sur \texttt{RisksPage} et \texttt{TeamPage} pour le scan Watchdog (\texttt{POST /webhook/api/watchdog/scan}), et dans \texttt{project-viability-refresh.ts} pour la viabilité.

Le Helper (\texttt{helper.api.ts}, \texttt{chat.api.ts}) envoie le contexte projet et la conversation à \texttt{POST /webhook/api/helper/chat} ; la réponse structurée inclut texte et \texttt{suggested\_actions}. Le panneau Copilot global peut appeler \texttt{POST /webhook/v1/copilot/insights} (\texttt{copilot-provider.tsx}) avec un contexte portfolio construit par \texttt{buildInsightsContextPayload}.

\subsection{Exposition des APIs via Webhooks}

\texttt{httpClient} (axios) attache le Bearer JWT sur chaque requête et gère le refresh sur 401 via \texttt{/webhook/refresh}. \texttt{apiGet} et \texttt{apiPost} (\texttt{apiClient.ts}) servent les routes \texttt{/api/copilot/*} et workspace talent. Les réponses n8n souvent imbriquées (\texttt{json}, \texttt{body}) sont aplaties par \texttt{flattenN8nRow} dans \texttt{use-reports-n8n.ts} et par des parseurs métier (\texttt{copilot-api.parse.ts}, \texttt{rh-api-parse.ts}).

Le tableau~\ref{tab:impl-api-ch5} rappelle les endpoints les plus sollicités lors de la réalisation.

\begin{table}[htbp]
\centering
\caption{Endpoints critiques de la réalisation}
\label{tab:impl-api-ch5}
\begin{tabular}{@{}ll@{}}
\toprule
\textbf{Endpoint} & \textbf{Fichier client principal} \\
\midrule
\texttt{POST /webhook/api/project/viability} & \texttt{orchestrator.api.ts} \\
\texttt{POST /webhook/api/copilot/recompute} & \texttt{recompute.api.ts}, \texttt{useProjects.ts} \\
\texttt{GET /webhook/api/copilot/dashboard} & \texttt{copilot.api.ts} \\
\texttt{GET /webhook/wmp-detail-v1/.../projects/:id} & \texttt{manager-projects.api.ts} \\
\texttt{POST /webhook/api/helper/chat} & \texttt{helper.api.ts} \\
\texttt{GET /api/workspace/talent/*} & \texttt{talent-workspace-pages.api.ts} \\
\bottomrule
\end{tabular}
\end{table}

\section{Implémentation de la base de données}

\subsection{Structure Supabase/PostgreSQL}

PostgreSQL héberge les tables métier et les tables de résultats analytiques. Supabase fournit l'hébergement managé et le bucket pour les PDF de rapports (\texttt{file\_url} renvoyé par \texttt{GenerateReportResponse}). Les workflows n8n accèdent aux tables via des nœuds Postgres ; le front-end n'exécute pas de SQL.

\subsection{Principales tables}

Les entités manipulées indirectement par l'UI correspondent aux types TypeScript et au cahier Agentic AI. \texttt{projects} porte les scores \texttt{viability\_score}, \texttt{risk\_score} et la décision en lecture seule côté client. \texttt{project\_tasks}, \texttt{project\_requirements} et \texttt{project\_assignments} structurent le travail et le staffing. \texttt{talents}, \texttt{skills} et \texttt{talent\_skills} modélisent le référentiel RH. \texttt{workload\_snapshots} alimente les vues de charge. \texttt{project\_viability\_scores} et \texttt{ai\_decisions\_log} historisent les sorties orchestrateur et les validations humaines.

Les demandes RH et alertes sont matérialisées dans des tables ou vues backend exposées sous forme de DTO (\texttt{RhActionItem}, notifications dans \texttt{notifications.api.ts}).

\subsection{Contraintes et intégrité des données}

L'intégrité repose sur des UUID validés côté client avant envoi (\texttt{assertUuid} dans \texttt{manager-api-contract}, \texttt{rh-actions.api.ts}). \texttt{orchestrator.api.ts} sanitise les identifiants talent pour le what-if afin d'éviter les erreurs serveur sur UUID invalides. Les corps POST d'affectation (\texttt{wmp-assign-v1}) imposent des champs nullables explicites conformément au contrat n8n.

Les mutations CRUD projet peuvent renvoyer un bloc \texttt{analysis\_refresh} (\texttt{CrudMutationResult}) pour afficher les scores à jour sans recalcul local. Les contraintes métier (statuts projet, types de demandes RH) sont appliquées dans les workflows ; le front-end reflète les énumérations reçues (\texttt{Continue}, \texttt{Adjust}, \texttt{Stop}, types \texttt{skill\_gap}, \texttt{training}, etc.).

\section{Interfaces principales}

\subsection{Tableau de bord Manager}

\texttt{DashboardPage} agrège les indicateurs via \texttt{useDashboard} (\texttt{GET /webhook/manager/dashboard} et données Copilot associées). La page affiche des KPI avec sparklines SVG, les sections \texttt{AnalystSection} et les recommandations Matchmaker construites par \texttt{buildExecutiveRecommendations} et \texttt{buildSkillGapsSorted} à partir des données serveur, sans inventer de scores. Des liens ouvrent Mission Control sur un projet via \texttt{managerProjectsOpenModalPath}.

La figure~\ref{fig:ui-dashboard-ch5} indique l'emplacement prévu pour une capture d'écran du tableau de bord (à insérer sous \texttt{rapport/figures/dashboard-manager.png} si disponible).

\begin{figure}[htbp]
\centering
\IfFileExists{figures/dashboard-manager.png}{%
  \includegraphics[width=0.92\textwidth]{figures/dashboard-manager.png}}{%
  \fbox{\parbox{0.9\textwidth}{\centering\vspace{1.2cm}
  \textit{Figure : tableau de bord manager}\\
  KPI, blocs Analyst et Matchmaker, liens vers projets\vspace{1.2cm}}}}
\caption{Tableau de bord de l'espace manager (\texttt{DashboardPage})}
\label{fig:ui-dashboard-ch5}
\end{figure}

\subsection{Espace RH}

Les pages RH consomment \texttt{rh-workspace.api.ts} : \texttt{fetchRhDashboard} sur \texttt{/api/rh/dashboard}, \texttt{fetchRhCriticalGaps}, \texttt{fetchRhTrainingPlans}, \texttt{fetchRhOrgAlerts}, ainsi que les endpoints de simulation de réallocation (\texttt{/api/rh/reallocation/simulate} et \texttt{validate}). La gestion des comptes utilise \texttt{listRhUsers}, \texttt{createRhUser}, \texttt{patchRhUserRole} et \texttt{patchRhUserStatus} vers les webhooks \texttt{/webhook/rh/users}. Le catalogue de compétences et la mobilité s'appuient sur les hooks CRUD \texttt{skills} et \texttt{talents} lorsque configurés.

\subsection{Espace Talent}

Chaque page talent est implémentée dans \texttt{talent-workspace-pages.tsx} et charge ses données via \texttt{use-talent-workspace-page-query} et le chemin dédié \texttt{/api/workspace/talent/...}. Le tableau de bord talent, la liste des projets assignés, les tâches, la charge, les compétences et les formations sont en lecture seule côté métier analytique ; seule la mise à jour de tâche autorise une mutation explicite. Les notifications peuvent utiliser \texttt{VITE\_TALENT\_NOTIFICATIONS\_URL} ou l'agrégat workspace.

\subsection{Mission Control}

\texttt{ProjectMissionControlModal} (\texttt{project-mission-control-modal.tsx}, plus de 1\,000 lignes) concentre le pilotage d'un projet. Six onglets sont implémentés : \texttt{overview}, \texttt{team}, \texttt{tasks}, \texttt{risks}, \texttt{simulation} et \texttt{decisions}, correspondant au type \texttt{MissionControlWorkspaceTabId}. Le modal charge le détail via \texttt{wmp-detail-v1}, affiche le stepper de cycle de vie, les traces d'agents (\texttt{buildAgentTraces}), permet assign/unassign, la gestion des tâches, le rafraîchissement de viabilité (\texttt{use-project-viability-refresh}) et l'intégration du panneau risques (\texttt{ProjectMissionControlRisks}).

\begin{figure}[htbp]
\centering
\fbox{\parbox{0.9\textwidth}{\centering\vspace{1cm}
\textit{Figure : Mission Control — onglets overview, équipe, tâches, risques, simulation, décisions}\vspace{1cm}}}
\caption{Modal \texttt{ProjectMissionControlModal} (centre de pilotage projet)}
\label{fig:ui-mission-ch5}
\end{figure}

\subsection{Copilot conversationnel}

Deux niveaux coexistent. Le \texttt{CopilotPanel} global (\texttt{CopilotProvider}) envoie le contexte de page (\texttt{dashboard}, \texttt{projects\_list}, \texttt{project\_detail}, \texttt{staffing}) vers les insights et permet d'enregistrer une décision stratégique (\texttt{saveCopilotDecision}, \texttt{postStrategicDecision}). La page \texttt{HelperChatPage} implémente le dialogue multi-tours avec historique (\texttt{managerConversationsApi}, archivage \texttt{PATCH .../archive}). \texttt{ManagerProjectCopilotPanel} restreint le Copilot au contexte d'un projet ouvert dans Mission Control.

\subsection{Simulateur What-if}

\texttt{WhatIfProvider} enveloppe l'application et expose \texttt{open()} / \texttt{close()} pour lancer une simulation depuis plusieurs écrans. \texttt{ProjectWhatIfSimulator} collecte les paramètres ; l'appel principal utilise \texttt{orchestratorApi.whatIf} (\texttt{POST /webhook/api/project/what-if}) ou \texttt{postCopilotProjectWhatIf} selon le chemin configuré. \texttt{WhatIfResultPanel} affiche score, décision, breakdown et recommandations renvoyés par le serveur, sans recalcul local.

\section{Fonctionnalités avancées}

\subsection{Alertes intelligentes}

\texttt{RisksPage} et \texttt{use-manager-risk-data} consomment les risques projet via POST ou GET sur \texttt{/webhook/api/project/risks}. Le scan Watchdog (\texttt{POST /webhook/api/watchdog/scan} avec \texttt{use\_ai: true}) peut être déclenché depuis l'équipe. Les notifications manager (\texttt{NotificationsPage}, barre supérieure) agrègent alertes risque et actions RH. Le PATCH \texttt{/webhook/manager/risk-alerts/:id} permet d'acquitter ou ignorer une alerte. Côté RH, \texttt{RhOrgAlertsPage} charge \texttt{/api/rh/organizational-alerts}.

\subsection{Recommandations IA}

Les recommandations proviennent des blocs \texttt{matchmaker} et \texttt{analyst} du dashboard, du matching projet (\texttt{talent-matching-panel}), du Copilot insights, du Helper (\texttt{suggested\_actions}) et des options Strategist. L'affichage utilise \texttt{formatUserFacingExplanation} et \texttt{stripTechnicalScoringSegments} pour présenter des textes métier lisibles. Le paramètre \texttt{use\_ai: true} est systématique sur les analyses à forte valeur ajoutée côté manager.

\subsection{Journal des décisions}

\texttt{DecisionLogPage} appelle \texttt{useDecisions} (\texttt{GET /webhook/manager/copilot-decisions}) et visualise la distribution Continue / Adjust / Stop avec Recharts (barres et courbes, couleurs sémantiques \texttt{\#059669}, \texttt{\#d97706}, \texttt{\#dc2626}). Chaque entrée peut renvoyer vers Mission Control ou les demandes RH. L'enregistrement d'une nouvelle décision passe par les API Copilot décrites en section~4.

\subsection{Génération de rapports}

\texttt{ReportsPage} et \texttt{use-reports-n8n.ts} implémentent le cycle complet : synthèse (\texttt{GET /webhook/reports/summary}), génération board-pack ou dossier projet (\texttt{POST /webhook/reports/generate-board-pack}, \texttt{generate-project-dossier}), historique (\texttt{GET /webhook/reports/history}), envoi courriel (\texttt{POST /webhook/reports/send-email}) et planification (\texttt{POST /webhook/reports/schedule}). Les PDF sont stockés côté Supabase ; l'UI affiche \texttt{file\_url} pour téléchargement. Les erreurs utilisent \texttt{skipGlobalHttpErrorToast} pour un traitement localisé.

\section{Difficultés rencontrées et solutions apportées}

La réalisation a confronté plusieurs difficultés techniques, documentées implicitement dans le code par des contournements explicites.

La \textbf{latence et la variabilité des réponses n8n} ont imposé des timeouts configurables (\texttt{VITE\_HTTP\_TIMEOUT\_MS}, \texttt{VITE\_PORTFOLIO\_TIMEOUT\_MS}) et l'usage de React Query pour ne pas bloquer l'UI pendant les analyses longues. Le proxy Vite utilise un délai de 60 secondes pour les mutations lourdes.

Le \textbf{format JSON encapsulé} des webhooks (champs \texttt{json} ou \texttt{body}) a nécessité des fonctions d'aplatissement (\texttt{flattenN8nRow}, \texttt{unwrapDataPayload}, parseurs RH et Copilot) pour stabiliser l'affichage indépendamment des variations de sortie des nœuds n8n.

Les \textbf{contraintes CORS en développement} ont été résolues par le proxy \texttt{/webhook}, \texttt{/api} et \texttt{/n8n-webhook} plutôt que par des appels directs cross-origin depuis le navigateur, sauf configuration explicite \texttt{VITE\_N8N\_DIRECT\_IN\_DEV=1}.

La \textbf{coexistence de plusieurs conventions d'URL} (chemins versionnés \texttt{wmp-*}, webhooks UUID, routes \texttt{/api/rh/actions} distinctes) est gérée par des fichiers de configuration (\texttt{manager-projects-api.config.ts}, \texttt{wmp-assignments-webhook.config.ts}, \texttt{crud-api.ts}) et par des variables \texttt{VITE\_*} documentées dans \texttt{vite-env.d.ts}.

La \textbf{gestion des sessions JWT} (expiration, refresh, appels \texttt{/webhook/auth/me} intempestifs) est centralisée dans \texttt{http-client.ts} et \texttt{auth-provider.tsx}, avec file d'attente des requêtes pendant le refresh.

Enfin, le respect strict de la règle « pas de calcul front » a parfois complexifié l'UX (attente des workflows après mutation) ; le recompute automatique optionnel après affectation et le bloc \texttt{analysis\_refresh} après CRUD compensent en rafraîchissant les scores affichés dès que le serveur répond.

\section{Conclusion}

Ce chapitre a présenté la réalisation et l'implémentation de Strategic Copilot : environnement de développement npm/Vite, stack React 19 et TypeScript, configuration proxy et production Vercel, structure du front-end et gestion d'état avec React Query, implémentation des workflows n8n et de l'intégration OpenAI côté serveur, modèle de données PostgreSQL/Supabase, interfaces manager, RH et talent (dont Mission Control, Copilot, what-if et rapports), ainsi que les difficultés techniques rencontrées et les solutions retenues.

Le livrable \texttt{copilot-ui} constitue une application SaaS opérationnelle, alignée sur le principe API-first et sur l'architecture multi-agents documentée dans les chapitres précédents. Les captures d'écran peuvent être ajoutées dans \texttt{rapport/figures/} pour enrichir les figures~\ref{fig:ui-dashboard-ch5} et~\ref{fig:ui-mission-ch5} lors de la version finale du mémoire. Le chapitre suivant présente la stratégie de test, les résultats des validations et l'évaluation globale de la solution sur l'environnement de préproduction.
